\section{Обсуждение}\label{sec:discussion}

\textbf{Что подтвердилось и что нет.} Запас, откалиброванный по одной ошибке, перестаёт работать, когда действует вторая (H1, H2), а совместный запас не больше суммы двух отдельных и при большом дрейфе меньше её на четверть (H3). Связи mIoU с безопасностью на 15 сценах не видно, но и отсутствие её не доказано (H4). Три ожидания не оправдались. Положительная корреляция ошибок делает сумму запасов консервативнее, а недопокрытие возникает, когда редкие сбои локализации и ошибки меток приходятся на разные сцены, и оно невелико. Множества меток выродились не из-за полных пропусков, а из-за обрезанных краёв. Реальная одометрия ошибается не плавно, как модель дрейфа, а скачками.

\textbf{Отрицательные результаты} (получены на предыдущей версии конвейера и не пересчитывались). Выпуклая оболочка и морфологическое замыкание областей не меняют промах: пропущенная часть объекта лежит вне найденного фрагмента, а не в дырах внутри него. Запас, зависящий от уверенности фрагмента, не уменьшил запретную зону.

\textbf{Главное узкое место} --- уверенный пропуск или подмена класса целого объекта. Пока калибровочных сцен не меньше $1/\alpha - 1$, но меньше $2/\alpha - 1$, запас равен наихудшему промаху, и одна такая сцена заставляет класс отказываться в значительной доле разбиений; при 19 сценах растение почти не отказывается. Более сильные детекторы убирают пропуск растения, но у каждого из трёх проверенных находится свой целиком пропущенный класс; подмену тумбы создаёт словарь, а закрытый словарь превращает её в ложные срабатывания. Объединение детекторов по максимуму это узкое место снимает, а проверка находок и голосование возвращают его: гарантия чувствительна к полноте, а не к точности. Метрики качества карты этого не отражают, и в протокол бенчмарков вроде OSMa-Bench естественно добавить долю сцен, где объект пропущен или подменён целиком.

\textbf{Какую гарантию требовать.} Гарантия «каждый объект внутри запаса» дорога, потому что за неё платят и объекты, мимо которых не идёт ни один путь. Гарантия на ожидаемую долю опасных миссий (разд.~\ref{sec:crc}) почти не уменьшает число решённых задач, но допускает целиком опасную сцену. Для реальной локализации важен хвост ошибок: SLAM с замыканиями втрое снижает медианную ошибку, но запас задают сцены без замыканий.

\textbf{Ограничения.}
\begin{itemize}[nosep]
  \item Одно семейство из 21 симулированной сцены (пять групп расстановок мебели) и одно условие освещения; бутстреп по сценам эту кластеризацию не учитывает. Покрытие на случайных разбиениях одних и тех же сцен выполняется по построению; о поведении на новых сценах говорят лишь бутстреп-интервалы шириной 0,1--0,15, проверка с отложенной планировкой и две сцены HM3D.
  \item Основной дрейф синтетический и не зависит от ошибок детектора. Реальная локализация проверена одометрией и собственным SLAM с замыканиями на основе Open3D; промышленные системы (ORB-SLAM3, RTAB-Map) не проверялись.
  \item Восприятие --- детектор рамок с масками по глубине или MobileSAM, а не полный картограф.
  \item Гарантия маргинальна по калибровочным выборкам и поклассова; для нескольких классов сразу она слабее, $1 - |\mathcal K_m|\,\alpha$.
  \item Основная гарантия относится к плану в момент запроса (допущение~\ref{as:exec}). Исполнение с дрейфом проверено моделью дрейфа, а не реальным роботом; запас на исполнение калибруется по объединённой выборке путей, и его обменяемость приближённая. Столкновения с препятствиями не покрываются. Дискретизация может съесть до 7~см безопасного расстояния.
  \item Семейство трудных задач и пара запретных классов для миссий выбраны после первых прогонов; множества меток реализованы нами на итоговой карте, без режима исследования авторов.
\end{itemize}

\section{Заключение}

Расстояние промаха в системе координат планировщика --- одна калибруемая величина, которая даёт планировщику гарантию сразу против ошибки метки и ошибки позы и честно отказывается, когда гарантия была бы бесполезной. Запасы по одной ошибке теряют покрытие под действием другой, сумма отдельных запасов переплачивает или недопокрывает в зависимости от неизвестного расположения сбоев, а калибровка множеств меток вырождается в запрет всего занятого. Цену гарантии задаёт наихудшая ошибка восприятия или локализации среди калибровочных сцен. Её снижают объединение детекторов (полнота важнее точности) и переход к гарантии на долю опасных миссий; исполнение с дрейфом до 7~см не ломает её относительно запретных классов, хотя 1--9\,\% исполнений задевают препятствия. Следующие шаги --- больше калибровочных сцен (например, данные OSMa-Bench++), SLAM, находящий замыкания во всех сценах, перелокализация во время движения и полный открытословарный картограф вместо детектора рамок.
